% lemniro: {"kind":"journal","description":"The small word finite does a large amount of work in topology. A shrinking-interval example explains why.","topic":"Topology","date":"2026-10-01","readingMinutes":5,"prerequisites":["Open intervals"],"related":["topological-spaces","compact-sets-are-closed"]}
\documentclass[11pt]{article}
\input{tex/lemniro-preamble}
\title{The work done by the word finite}
\begin{document}
\maketitle

\section{An intersection that loses its room}
For each positive integer $n$, let $U_n=(-1/n,1/n)$. Every $U_n$ gives
zero some room: it contains a whole interval around zero. Any finite
selection of these intervals still does so. If the largest selected
index is $N$, then the intersection of that selection is exactly
$U_N$.

The infinite intersection behaves differently:
\begin{equation}\label{eq:intersection}
  \bigcap_{n=1}^{\infty}U_n=\set{0}.
\end{equation}
Zero belongs to every interval, so it belongs to the intersection. If
$x\ne0$, choose an integer $n>1/|x|$. Then $1/n<|x|$, which means that
$x\notin U_n$. Thus no nonzero point belongs to the intersection,
proving Equation~\ref{eq:intersection}.

The resulting singleton is not open in $\R$. Each finite stage leaves
an open neighborhood, while the intersection of all stages does not.
There is no contradiction: arbitrary intersections of open sets are
not among the topology axioms.

\section{A proof-reading habit}
Whenever a proof intersects open neighborhoods, pause at the index set.
If the author has just chosen a finite subcover, that choice may be
doing more than shortening notation. It may be the exact reason the
next intersection remains open.

In the compact-Hausdorff argument, we begin with one separating
neighborhood for each point of a compact set. Compactness lets us keep
only finitely many. Their intersection then supplies a neighborhood
of the point outside the set. The word ``finite'' is the bridge from
a family of local choices to a valid open set.

\begin{remark}
An infinite intersection can still be open. For example,
$\bigcap_{n\in\N}\R=\R$. Equation~\ref{eq:intersection} shows that
openness is not guaranteed; it does not say that infinite
intersections are always nonopen.
\end{remark}
\end{document}
